\providecommand{\CY}{\mathrm{CY}}
\providecommand{\kch}{\kappa_{\mathrm{ch}}}
\providecommand{\DT}{\mathrm{DT}}
\providecommand{\PT}{\mathrm{PT}}
\providecommand{\NCDT}{\mathrm{NCDT}}
\providecommand{\Ytil}{\widetilde Y}
\providecommand{\Yplus}{Y^{+}}
\providecommand{\MacM}{M}
\providecommand{\chtop}{\chi_{\mathrm{top}}}
\providecommand{\CoHA}{\mathrm{CoHA}}
\providecommand{\Jac}{\mathrm{Jac}}
\providecommand{\bP}{\mathbb{P}}
\providecommand{\C}{\mathbb{C}}
\providecommand{\gr}{\mathrm{gr}}

\section*{Conifold Hall character and MacMahon specialization}

\subsection*{Resolved conifold datum}

Let
\[
  X_{\mathrm{con}}
  =
  \Ytil_{\mathrm{con}}
  =
  \operatorname{Tot}\bigl(\mathcal O_{\bP^1}(-1)
  \oplus \mathcal O_{\bP^1}(-1)\bigr)
\]
be a small resolution of the ordinary conifold
\(\{xy-zw=0\}\subset \mathbb A^4\).  The compact curve is the zero
section \(\bP^1\), denoted by \(\beta\), with normal bundle
\(\mathcal O(-1)\oplus\mathcal O(-1)\).  This is a noncompact toric
\(\CY_3\) of curve-local type.  It lies outside the local-surface class
\(\operatorname{Tot}(K_S\to S)\), whose base is a surface and whose
canonical-bundle fibre has rank one.

Use the signed Donaldson--Thomas variable
\[
  s=-q_{\DT}, \qquad
  \MacM(s)=\prod_{m\geq 1}(1-s^m)^{-m},
\]
and let \(Q\) record the curve class \(\beta\).  The large-volume
Gopakumar--Vafa input is the primitive rational curve
\[
  n^0_{\beta}=1,\qquad
  n^g_{d\beta}=0 \quad\text{for }(g,d)\ne(0,1).
\]
Every nonzero \(Q^d\)-coefficient with \(d>1\) in the logarithm is a
multi-cover contribution of this primitive class.  Consequently the
reduced stable-pairs factor is
\[
  Z^{\PT}_{\mathrm{con}}(s,Q)
  =
  \prod_{m\geq 1}(1-Qs^m)^m .
\]

\subsection*{Large-volume Hall character}

Let \(\Yplus_{\theta}(X_{\mathrm{con}})\) denote the positive Hall object
in the large-volume stability chamber \(\theta\).  Its scalar graded
Euler series is the rank-one DT series
\[
  \chi_{\gr}\!\left(\Yplus_{\theta}(X_{\mathrm{con}})\right)(s,Q)
  =
  Z^{\DT}_{\mathrm{con}}(s,Q)
  =
  \MacM(s)^2 \prod_{m\geq 1}(1-Qs^m)^m .
  \tag{1}
\]
\(\MacM(s)^2\) is the degree-zero Hilbert-scheme factor
\(\MacM(s)^{\chtop(X_{\mathrm{con}})}\), since
\(\chtop(X_{\mathrm{con}})=\chi(\bP^1)=2\).  The second factor is forced
by the single primitive class \(n^0_{\beta}=1\).

The equality in \((1)\) is a character identity.  The algebra object is
the critical CoHA of the Klebanov--Witten Jacobi algebra
\[
  \CoHA\!\left(\Jac(Q_{\mathrm{con}},W_{\mathrm{con}})\right),
  \qquad
  W_{\mathrm{con}}
  =
  \operatorname{tr}(a_1b_1a_2b_2-a_1b_2a_2b_1).
\]
With the Li--Yamazaki current presentation, the Davison--Meinhardt PBW
theorem, and the Morrison--Mozgovoy--Nagao--Szendr\H{o}i root
decomposition as input, this positive half is the conifold current
algebra \(Y^+(\widehat{\mathfrak{gl}}(1|1))^{\mathrm{con}}\) as a
\(\mathbb Z_{\geq 0}^2\)-graded connected bialgebra.  This is a
positive-half statement.  The full Yangian, \(R\)-matrix, and
vertex-algebra realization require Drinfeld double, centre, completion,
or evaluation operations.  At \(d=3\), the direct \(\Phi_3\)-output is
\(E_1\)-chiral; \(E_2\)-structure belongs to the centre or double stage.

The MacMahon factor in \((1)\) is degree-zero DT data.  It is separate
from the \(\mathbb C^3\) identity
\(\CoHA(\mathbb C^3)=Y^+(\widehat{\mathfrak{gl}}_1)\).  The conifold
CoHA is the Klebanov--Witten/NCCR Hall algebra, and a comparison with
the \(\mathbb C^3\) positive half requires the explicit wall-crossing
or chart-gluing datum.

\subsection*{Szendr\H{o}i chamber and the \(Q^{-1}\)-factor}

The noncommutative chamber attached to the Klebanov--Witten NCCR keeps
both orientations of the exceptional curve.  Its two-variable scalar
product is
\[
  Z^{\NCDT}_{\mathrm{con}}(s,Q)
  =
  \MacM(s)^2
  \prod_{m\geq 1}(1-Qs^m)^m
  \prod_{m\geq 1}(1-Q^{-1}s^m)^m .
  \tag{2}
\]
The \(Q^{-1}\)-factor is the flopped curve contribution.  Passing from
\((2)\) to the large-volume series \((1)\) is the Nagao--Nakajima
wall-crossing; it removes the flopped factor by changing stability
chamber.  The \(Q\)-grading is part of the character datum.
The specialization \(Q=1\) of the noncommutative product forgets the two
curve orientations and leaves the trivial scalar
\[
  \left. Z^{\NCDT}_{\mathrm{con}}(s,Q)\right|_{Q=1}
  =
  \MacM(s)^2
  \left(\prod_{m\geq 1}(1-s^m)^m\right)^2
  =
  1.
\]
This scalar differs from the large-volume specialization in \((3)\).
Chamber comparison must be made before \(Q\) is forgotten.

\subsection*{One-variable large-volume specialization}

After the wall-crossing to the large-volume chamber, setting \(Q=1\) in
\((1)\) gives
\[
  \chi_{\gr}\!\left(\Yplus_{\theta}(X_{\mathrm{con}})\right)(s,1)
  =
  \MacM(s)^2\prod_{m\geq 1}(1-s^m)^m
  =
  \MacM(s).
  \tag{3}
\]
The one-variable scalar character is
\[
  \boxed{
  \chi_{\gr}\!\left(\Yplus_{\theta}(X_{\mathrm{con}})\right)(s,1)
  =
  1+s+3s^2+6s^3+13s^4+24s^5+48s^6+86s^7+\cdots .}
\]
The equality with \(\MacM(s)\) is a scalar equality after the curve
class has been forgotten.  The conifold Klebanov--Witten CoHA and the
\(\mathbb C^3\) positive half remain distinct Hall algebras.

The first coefficients follow from
\[
  \MacM(s)^2
  =
  1+2s+7s^2+18s^3+47s^4+110s^5+258s^6+\cdots
\]
together with
\[
  \prod_{m\geq 1}(1-s^m)^m
  =
  \MacM(s)^{-1}
  =
  1-s-2s^2-s^3+4s^5+4s^6+7s^7+\cdots .
\]
Convolution gives
\[
\begin{aligned}
  [s^0]&=1,\\
  [s^1]&=-1+2=1,\\
  [s^2]&=-2-2+7=3,\\
  [s^3]&=-1-4-7+18=6,\\
  [s^4]&=0-2-14-18+47=13.
\end{aligned}
\]

\subsection*{The conifold modular characteristic}

The chiral modular characteristic of the conifold row is
\[
  \kch(A_{X_{\mathrm{con}}})=1.
\]
This is the direct McKay/NCCR calculation for the Klebanov--Witten
two-node conifold algebra.  The Jacobi algebra is module-finite over
the conifold centre and has crepant global dimension three; its two
small-resolution chambers carry the same single compact curve class.
Equivalently, the primitive GV input is \(n^0_{\beta}=1\), giving one
primitive Hall mode.  The number \(1\) is independent of the degree-zero
DT exponent \(2=\chtop(X_{\mathrm{con}})\), the chamber variable \(Q\),
and any Borcherds/BKM weight.  The compact Hodge-supertrace formula is
unavailable on this noncompact \(\CY_3\); the conifold row is computed
from the local Jacobi algebra and its Hall primitive.

\subsection*{Verification}

\emph{Product arithmetic.}  The products
\(\MacM(s)^2\), \(\MacM(s)^{-1}\), and their convolution give
\((3)\) coefficient by coefficient after \(Q\) is forgotten in the
large-volume chamber.

\emph{GV/PT route.}  The single primitive invariant
\(n^0_{\beta}=1\) gives
\[
  Z^{\PT}_{\mathrm{con}}(s,Q)
  =
  \prod_{m\geq 1}(1-Qs^m)^m ,
\]
and the DT/PT factorization contributes the degree-zero term
\(\MacM(s)^2\).

\emph{CoHA route.}  The Klebanov--Witten critical CoHA gives the
positive Hall object.  Its chamber character is the DT product in
\((1)\).  Multiplication, coproduct, PBW root decomposition, and the
current presentation are additional structure beyond the scalar
character.

\emph{Local compute route.}  The local engine
\texttt{compute/lib/toric\_cy3\_dt\_engine.py} implements
\[
  \texttt{conifold\_reduced}=\prod_{m\geq 1}(1-Qs^m)^m,\qquad
  \texttt{conifold\_full}=\MacM(s)^2\texttt{conifold\_reduced}.
\]
The tests in \texttt{compute/tests/test\_toric\_cy3\_dt\_engine.py}
verify the dual Cauchy expansion, the extraction
\(n^0_{\beta}=1\), the vanishing \(n^0_{d\beta}=0\) for \(d>1\), and the
logarithmic identity
\[
  [Q^d]\log Z^{\PT}_{\mathrm{con}}
  =
  -\frac{1}{d}\frac{s^d}{(1-s^d)^2}.
\]

\subsection*{References}

Szendr\H{o}i, B. \emph{Non-commutative Donaldson--Thomas invariants and
the conifold.} Geom. Topol. 12 (2008), 1171--1202.

Morrison, A.; Mozgovoy, S.; Nagao, K.; Szendr\H{o}i, B.
\emph{Motivic Donaldson--Thomas invariants of the conifold and the
refined topological vertex.} Adv. Math. 230 (2012), 2065--2093.

Nagao, K.; Nakajima, H. \emph{Counting invariant of perverse coherent
sheaves and its wall-crossing.} Int. Math. Res. Not. (2011), 3885--3938.

Maulik, D.; Nekrasov, N.; Okounkov, A.; Pandharipande, R.
\emph{Gromov--Witten theory and Donaldson--Thomas theory, I.}
Compos. Math. 142 (2006), 1263--1285.

Pandharipande, R.; Thomas, R. \emph{Curve counting via stable pairs in
the derived category.} Invent. Math. 178 (2009), 407--447.

Davison, B.; Meinhardt, S. \emph{Cohomological Donaldson--Thomas theory
of a quiver with potential and quantum enveloping algebras.} Invent.
Math. 221 (2020), 777--871.

Li, W.; Yamazaki, M. \emph{Quiver Yangian from crystal melting.}
JHEP 11 (2020), 035.

Van den Bergh, M. \emph{Non-commutative crepant resolutions.}
The legacy of Niels Henrik Abel, Springer, 2004, 749--770.

MacMahon, P. A. \emph{Combinatory Analysis}, Vol. II. Cambridge
University Press, 1916.
